
Validation results are presented for three research questions. RQ1 establishes perfect linear correlation (r=1.000, p<0.0001) between micro-pilot and full-scale overhead. RQ2 shows Bayesian updates reduce prediction error by 40.91\% (p=0.0003). RQ3 validates Gate 1 classification accuracy at 93.3\% (p=4.$34\times 10^{-7})$, exceeding the 80\% target by 13.3 percentage points.

\subsubsection{5.1 RQ1: Micro-Pilot Correlation (Mechanism M1)}

Micro-pilot overhead $O_{10}$ exhibited perfect correlation with full-scale overhead O\_full across all 32 hypotheses.

\begin{table}[htbp]
\centering
\resizebox{\columnwidth}{!}{%
\begin{tabular}{llll}
\toprule
Metric & Result & Success Criterion & Status \\
\midrule
Pearson r & 1.000 & r > 0.7 & PASS \\
p-value & < 0.0001 & p < 0.05 & PASS \\
$R^2$ & 1.000 & $R^2$ > 0.5 & PASS \\
Scaling factor k & 1.000 $\pm$ 0.000 & Report mean & 1.000 \\
CV across types & 0.00\% & CV < 30\% & PASS \\
\bottomrule
\end{tabular}
}
\end{table}

\textbf{Observations:}

Perfect linearity: All 32 hypotheses fall exactly on the regression line O\_full = 1.000 $\times O_{10}$. This perfect fit ($R^{2}=1.000)$ reflects synthetic corpus design, where overhead was generated deterministically with k=1.000.

Consistent scaling across types: All hypothesis types (attention, gradient, regularization, normalization) exhibit identical scaling factor k=1.000 with zero variance (CV=0.00\%). This validates Assumption A3 (k generalizes across types) in the idealized case but represents a synthetic artifact unlikely to hold in real data.

Implications: Perfect correlation enables zero-error extrapolation from 10-sample micro-pilot to full-scale prediction. Real-world validation will test whether r $\geq$ 0.7 threshold holds under measurement noise (hardware variance, I/O overhead, memory bottlenecks).

This result validates the mechanism logic of linear overhead scaling---that extrapolation from micro-pilot to full-scale is mathematically sound when correlation is strong. However, perfect r=1.000 is a synthetic artifact. Real data is expected to yield r=0.7-0.9 (still above threshold but not perfect).

\subsubsection{5.2 RQ2: Bayesian Error Reduction (Mechanism M2)}

Bayesian updates combining Gate 1 prior with Gate 2 likelihood reduced prediction error by 40.91\% on average across 20 hypotheses.

\begin{table}[htbp]
\centering
\resizebox{\columnwidth}{!}{%
\begin{tabular}{llll}
\toprule
Metric & Result & Success Criterion & Status \\
\midrule
Mean error reduction & 40.91\% & > 40\% & PASS \\
Paired t-test & t=4.453, p=0.0003 & p < 0.05 & PASS \\
Sample size & 20 hypotheses & $\geq$ 10 & PASS \\
Gate 1 mean error & 0.6966 & Report & Baseline \\
Gate 2 mean error & 0.1111 & Report & - \\
\bottomrule
\end{tabular}
}
\end{table}

\textbf{Observations:}

Marginal excess: Error reduction 40.91\% vs 40\% threshold represents 0.91 percentage point margin. Success criterion met but without large safety buffer. Paired t-test highly significant (p=0.0003 << 0.05), confirming result not due to chance.

Large absolute reduction: Gate 1 mean error 0.6966 $\rightarrow$ Gate 2 mean error 0.1111 represents 84\% absolute reduction.

Strong prior limits update room: With perfect scaling (k=1.000, r=1.000), Gate 1 prior already accurate. Bayesian update has limited refinement potential when prior variance is low. Real data with weaker correlation (r=0.7-0.8) is expected to show larger error reduction (>50\%) as higher prior uncertainty provides more room for likelihood-driven updates.

Bayesian mechanism validated: posterior refinement reduces prediction error as designed. Marginal result (40.91\% vs 40\%) reflects synthetic data artifact (strong prior). Real-world scenarios with measurement noise will likely show larger Bayesian benefit.

\subsubsection{5.3 RQ3: Gate 1 Classification Accuracy (Mechanism M3)}

Gate 1 viability classification achieved 93.3\% accuracy (28 of 30 correct predictions), significantly outperforming 50\% random baseline.

\begin{table}[htbp]
\centering
\resizebox{\columnwidth}{!}{%
\begin{tabular}{llll}
\toprule
Metric & Result & Success Criterion & Status \\
\midrule
Accuracy & 93.3\% (28/30) & > 80\% (24/30) & PASS \\
Binomial test & p=4.$34\times 10^{-7}$ & p < 0.05 & PASS \\
Performance gain & 43.3pp vs random & Report & - \\
\bottomrule
\end{tabular}
}
\end{table}

\textbf{Confusion Matrix:}

\begin{table}[htbp]
\centering
\resizebox{\columnwidth}{!}{%
\begin{tabular}{lll}
\toprule
 & Predicted Non-Viable & Predicted Viable \\
\midrule
Actual Non-Viable (26) & TP = 25 & FN = 1 \\
Actual Viable (4) & FP = 1 & TN = 3 \\
\bottomrule
\end{tabular}
}
\end{table}

\textbf{Key Metrics:}
\begin{itemize}
\item Recall on non-viable: 96.2\% (25/26 correct)
\item Recall on viable: 75.0\% (3/4 correct)
\end{itemize}

\textbf{Observations:}

Strong accuracy: 93.3\% vs 80\% target represents 13.3 percentage point margin. Binomial test p=4.$34\times 10^{-7}$ << 0.05 shows result highly significant against 50\% null hypothesis (random guessing). Performance gain: 43.3 percentage points over random baseline.

High recall: 96.2\% recall on non-viable hypotheses (25/26 identified) means only 1 false negative. This aligns with framework design priority: minimize missed non-viable hypotheses (which waste full implementation effort) at cost of occasional false positives (viable hypotheses incorrectly stopped).

Imbalanced corpus: 26 non-viable vs 4 viable (87\% non-viable prevalence) reflects strict 10\% threshold. Accuracy metric dominated by high recall on large non-viable class. False positive rate 25\% (1/4 viable stopped) less representative due to small viable sample. Balanced validation (50-50 prevalence) needed to test whether accuracy $\geq 80$\% holds with equal class distribution.

Error analysis: The single false negative (non-viable hypothesis predicted viable) occurred at 10.3\% overhead---borderline case just above 10\% threshold. The single false positive (viable hypothesis predicted non-viable) occurred at 9.7\% overhead---also borderline. Both errors fall within $\pm 0.5$\% of threshold, suggesting measurement noise or conservative prediction.

Gate 1 classification mechanism validated: 93.3\% accuracy significantly exceeds 80\% target and 50\% random baseline. High recall (96.2\%) demonstrates framework reliably identifies non-viable hypotheses at micro-pilot stage. Imbalanced corpus (87\% non-viable) may inflate accuracy metric; balanced validation recommended as future work.

\subsubsection{5.4 Summary}

All three mechanisms (M1, M2, M3) achieved success criteria with statistical significance:
\begin{itemize}
\item M1 (Scaling): r=1.000 > 0.7, p<0.0001
\item M2 (Bayesian): Error reduction 40.91\% > 40\%, p=0.0003
\item M3 (Classification): Accuracy 93.3\% > 80\%, p=4.$34\times 10^{-7}$
\end{itemize}

Perfect linearity (r=1.000, k=1.000, CV=0.00\%) is unlikely in real experiments. Expected real-world r=0.7-0.9, k variance 10-30\%, measurement noise reducing perfect correlation. Framework mechanics validated; external validity requires real corpus validation.
